\medskip

On munit le plan d'un repère orthonormé.

Pour tout entier naturel $n$, on considère la fonction $f_n$ définie sur $[0~;~ +\infty[$ par :

\[f_0(x) = \e^{- x}\quad \text{et, pour }\: n \geqslant 1, \: \: f_n(x) = x^n\e^{-x}.\]

Pour tout entier naturel $n$, on note $\mathcal{C}_n$ la courbe représentative de la fonction $f_n$.

\medskip

\emph{Les parties {\rm A} et {\rm B} sont indépendantes.}

\medskip

\textbf{Partie A : Étude des fonctions \boldmath $f_n$ pour $n \geqslant 1$\unboldmath}

\medskip

On considère un entier naturel $n \geqslant 1$.

\medskip

\begin{enumerate}
\item 
	\begin{enumerate}
		\item On admet que la fonction $f_n$ est dérivable sur $[0~;~+\infty[$.
		
Montrer que pour tout $x \geqslant 0$,

\[f'_n(x) = (n - x)x^{n-1}\e^{-x}.\]

		\item Justifier tous les éléments du tableau ci-dessous:

\medskip

\begin{center}
\psset{unit=1cm,arrowsize=2pt 3}
\begin{pspicture}(7,3.25)
\psframe(7,3.25)\psline(0,2.25)(7,2.25)\psline(0,2.75)(7,2.75)
\psline(1,0)(1,3.25)
\uput[u](0.5,2.65){$x$}\uput[u](1.2,2.65){$0$}\uput[u](4,2.65){$n$}\uput[u](6.6,2.65){$+ \infty$}
\rput(0.5,2.5){$f'_n(x)$}\rput[u](2.5,2.5){+}\rput(4,2.5){0}\rput(5.5,2.5){$-$}
\uput[u](1.2,0){0}\uput[u](6.8,0){0}\uput[d](4,2.25){$\left(\frac{n}{\text{e}}\right)^n$}
\psline{->}(1.5,0.4)(3.4,1.8)\psline{->}(4.4,1.8)(6.4,0.4)
\rput(0.5,1.125){$f_n$}
\end{pspicture}
\end{center}
	\end{enumerate}
\item Justifier par le calcul que le point A$\left(1~;~\e^{-1}\right)$ appartient à la courbe $\mathcal{C}_n$.
\end{enumerate}

\bigskip

\textbf{Partie B : Étude des intégrales \boldmath $\displaystyle\int_0^1 f_n(x)\,\text{d}x$\unboldmath{} pour \boldmath $n \geqslant 0$\unboldmath}

\medskip

Dans cette partie, on étudie les fonctions $f_n$ sur [0~;1] et on considère la suite $(I_n)$ définie pour tout entier naturel $n$ par:

\[I_n = \displaystyle\int_0^1 f_n(x)\,\text{d} x = \displaystyle\int_0^1 x^n \e^{-x}\,\text{d}x.\]

\medskip

\begin{enumerate}
\item Sur le graphique ci-dessous, on a représenté les courbes $\mathcal{C}_0, \mathcal{C}_1, \mathcal{C}_2, \mathcal{C}_{10}$\: et \: $\mathcal{C}_{100}$.

\psset{unit=10cm,arrowsize=2pt 3,comma=true}
\begin{pspicture}(-0.1,-0.05)(1.05,1.1)
\psaxes[linewidth=1.25pt,Dx=0.1,Dy=0.1]{->}(0,0)(0,0)(1.05,1.1)
\psplot[plotpoints=2000,linewidth=1.25pt,linecolor=red]{0}{1}{2.71828 x neg exp}\uput[u](0.75,0.47){\red $\mathcal{C}_0$}
\psplot[plotpoints=2000,linewidth=1.25pt,linecolor=blue]{0}{1}{2.71828 x neg exp x mul}\uput[u](0.75,0.35){\blue $\mathcal{C}_1$}
\psplot[plotpoints=2000,linewidth=1.25pt,linecolor=gray]{0}{1}{2.71828 x neg exp x  dup mul mul}\uput[u](0.75,0.27){\gray $\mathcal{C}_2$}
\psplot[plotpoints=2000,linewidth=1.25pt,linecolor=magenta]{0}{1}{2.71828 x neg exp x 10 exp mul}\uput[u](0.75,0.06){\magenta $\mathcal{C}_{10}$}
\psplot[plotpoints=3000,linewidth=1.25pt,linecolor=green]{0}{1}{2.71828 x neg exp x 100 exp mul}\uput[u](0.92,0.02){\green $\mathcal{C}_{100}$}
\end{pspicture}
	\begin{enumerate}
		\item Donner une interprétation graphique de $I_0$.
		\item Par lecture de ce graphique, quelle conjecture peut-on émettre sur la limite
de la suite $(I_n)$ ?
	\end{enumerate}
\item Calculer $I_0$.
\item 
	\begin{enumerate}
		\item Soit $n$ un entier naturel.
		
Démontrer que pour tout $x \in [0~;~1]$,\:

\[0 \leqslant x^{n+1} \leqslant x^n.\]

		\item En déduire que pour tout entier naturel $n$, on a :

\[0 \leqslant I_{n+1} \leqslant I_n.\]
	\end{enumerate}
\item Démontrer que la suite $(I_n)$ est convergente, vers une limite positive ou nulle que l'on notera $\ell$.
\item En utilisant une intégration par parties, démontrer que pour tout entier naturel $n$ on a :

\[I_{n+1} = (n + 1)I_n - \dfrac{1}{\e}.\]
\item
	\begin{enumerate}
		\item Démontrer que si $\ell > 0$, l'égalité de la question 5 conduit à une contradiction.
		\item Démontrer que $\ell = 0$. On pourra utiliser la question 6. a.

\medskip

On donne ci-dessous le script de la fonction mystere, écrite en langage Python.

On a importé la constante e.

\begin{center}
\begin{Verbatim}[frame=single]
def mystere(n):
  I = 1 - 1/e
  L = [I]
  for i in range(n):
     I = (i + 1)*I - 1/e
     L.append(I) 
  return L
\end{Verbatim}
\end{center}
	\end{enumerate}
\item Que renvoie mystere(100) dans le contexte de l'exercice ?

\end{enumerate}

\bigskip
